\begin{lemma}
\label{editorial-source-lemma-base-change-Jacobson}
Let $k$ be a field. Let $S$ be a $k$-algebra.
For any field extension $K/k$ whose cardinality is larger
than the cardinality of $S$ we have
\begin{enumerate}
\item for every maximal ideal $\mathfrak m$ of $S_K$ the field
$\kappa(\mathfrak m)$ is algebraic over $K$, and
\item $S_K$ is a Jacobson ring.
\end{enumerate}
It suffices more generally that $S$ have a $k$-algebra generating
set of cardinality less than $\#K$. If this set is finite, no
cardinality restriction on $K$ is needed.
\end{lemma}

\begin{proof}
Fix the given field extension $K/k$ and write
$S_K=K\otimes_kS$. If $(s_i)_{i\in I}$ generates $S$ as a
$k$-algebra, then $(1\otimes s_i)_{i\in I}$ generates $S_K$ as a
$K$-algebra. Indeed every tensor is a finite sum of terms
$\lambda\otimes s$, and an expression for $s$ as a polynomial
in finitely many $s_i$ gives the corresponding expression for
$\lambda\otimes s$, with all coefficients carried by $k\to K$.
Theorem \ref{editorial-source-theorem-uncountable-nullstellensatz} therefore applies
whenever $\#I<\#K$, and its finite-generator case applies without
that restriction. Taking the generating family indexed by every
element of $S$ gives the hypotheses and conclusion stated originally.
\end{proof}